\begin{corolario}[Cambio de ley hasta la primera salida]\label{macro:cambio}
La conclusión de la Proposición~\ref{macro:parada} vale bajo la ley $\nu$ con otras constantes
$C_{\delta,\varrho_{\rm mac}},c_{\delta,\varrho_{\rm mac}}>0$. En particular, con probabilidad al menos
$1-C_{\delta,\varrho_{\rm mac}}\exp(-c_{\delta,\varrho_{\rm mac}}n)$, en todos los tableros hasta $k=M$,
incluido ese
extremo, el punto de
Bethe es el máximo único y cumple cotas macroscópicas uniformes de difusión, vacancia y brecha.
\end{corolario}
\begin{proof}
Por \eqref{macro:predictor}, eligiendo los umbrales suficientemente pequeños se pueden fijar, por
ejemplo, una cota superior $4C_\nu/\kappa_\nu^2$ para $k\max_e p_e$, una cota inferior $\kappa_\nu/4$
para las vacancias y una cota inferior estrictamente menor que $c_\nu/8$ para la brecha. La
comparación de pesos en el Lema~\ref{macro:brecha} y la Proposición~\ref{macro:reequilibrado}
justifica esta última elección. Éstas son las cotas macroscópicas que se verifican en cada padre
regular del tramo macroscópico.

Para dicho padre y cualquiera de sus hijos, la expansión de un paso de \cite[Teorema~7.1 y
Corolario~7.2, con sus pruebas]{pred} da
\[
 L(B+e)-L(B)=\log[p_e(1+d_e)]+r_e,\qquad
 |d_e|\le C'_\delta/k,\qquad |r_e|\le C''_\delta/k^2,\qquad
 \sum_{e\in i}p_ed_e=0
\]
en cada fila disponible $i$. Aquí $d_e$ es la corrección relativa $\delta_c$
de~\cite[Teorema~7.1]{pred} y $r_e$ es el resto $r_c$ de su prueba. La construcción del hijo por
eliminación en \cite[Sección~6.5]{pred} permite usar la expansión también en el tablero de primera
salida. Sólo
se exige el control del padre. Las dos leyes usan la misma elección uniforme de la próxima fila.
Condicionalmente a esa fila, el cociente de la probabilidad de columna bajo $\nu$ y su probabilidad
$p_e$ bajo $P^{\rm B}$ es
\[
 \frac{(1+d_e)e^{r_e}}{\sum_{e'\in i}p_{e'}(1+d_{e'})e^{r_{e'}}} .
\]
Para $k\ge2C'_\delta$ todos los factores $1+d_{e'}$ son positivos. La cancelación por fila implica
$\sum_{e'\in i}p_{e'}(1+d_{e'})=1$, de modo que el denominador pertenece a
$[e^{-C''_\delta/k^2},e^{C''_\delta/k^2}]$. El módulo del logaritmo del cociente es a lo sumo
$2C'_\delta/k+2C''_\delta/k^2$.

A lo largo de una trayectoria parada en la primera salida o al llegar a $M$, el módulo del logaritmo
de la densidad relativa se acota por \[
\sum_{k=M+1}^{n}\left(\frac{2C'_\delta}{k}+\frac{2C''_\delta}{k^2}\right) \le
2C'_\delta\log(n/M)+2C''_\delta/M=O_\delta(1). \] La cota incluye la transición que produce la
salida. Al sumar sobre las historias paradas en las que la discrepancia alcanza el radio
$\varrho_{\rm mac}$ obtenemos, por la Proposición~\ref{macro:parada}, la cola exponencial afirmada.
Fuera de ese suceso se aplica la reconstrucción determinista en todos los tableros hasta $M$.
\end{proof}
